\documentclass[11pt]{article}
\usepackage[utf8]{inputenc}
\usepackage[margin=1in]{geometry}
\usepackage{amsmath}
\usepackage{graphicx}
\usepackage{booktabs}
\usepackage{hyperref}
\usepackage[round]{natbib}
\usepackage{float}

\title{Parafoveal Semantic Integration in Natural Reading Measured with Rapid Invisible Frequency Tagging and MEG}
\author{Author A$^{1}$, Author B$^{1}$, Author C$^{1,2}$\\
\small $^{1}$Centre for Human Brain Health, University of Birmingham, UK\\
\small $^{2}$School of Psychology, University of Birmingham, UK}
\date{}

\begin{document}
\maketitle

\begin{abstract}
Whether readers integrate semantic information from parafoveal words before fixating them remains debated, particularly in English. Existing paradigms such as the gaze-contingent boundary technique cannot distinguish pre-fixation integration from cross-saccade mismatch effects. Here, we combined Rapid Invisible Frequency Tagging (RIFT) at 60~Hz with MEG and eye tracking to provide a direct neural measure of parafoveal semantic processing during natural sentence reading. Thirty-four native English speakers read 160 sentences containing target nouns that were either congruent or incongruent with the sentence context. The target word region was flickered at 60~Hz using a 1440~Hz projector. In 29 participants with significant RIFT responses (identified via cluster-based permutation, $p < 0.01$), we found significantly weaker 60~Hz coherence for contextually incongruent versus congruent target words during the pre-target fixation ($p = 0.011$, cluster-based permutation). This congruency effect emerged within 100~ms of fixating the pre-target word. The RIFT congruency effect correlated with individual reading speed (Spearman $\rho = -0.39$, $p = 0.037$), linking parafoveal semantic integration to reading efficiency. RIFT source localised to left visual association cortex (Brodmann area 18, MNI $[-9, -97, 3]$). No congruency effect was observed during direct fixation of the target word. These results provide direct neural evidence that semantic information from parafoveal words is integrated with sentence context before fixation, and that the degree of integration predicts reading speed.
\end{abstract}

\section{Introduction}

Skilled reading depends on the ability to extract information not only from the currently fixated word but also from upcoming words in the parafovea, typically 2--5 degrees from fixation \citep{rayner1975parafoveal, rayner1998eye}. Parafoveal preview benefit---the facilitation of processing a word when it was available in the parafovea during the prior fixation---is well established for orthographic and phonological features \citep{schotter2012parafoveal}. Whether this benefit extends to semantic information, particularly in non-logographic languages such as English, has been intensely debated \citep{schotter2015does, hohenstein2010semantic, pan2021parafoveal}.

The primary tool for studying parafoveal processing has been the gaze-contingent boundary paradigm, in which the parafoveal preview is changed during the saccade to the target word \citep{rayner1975parafoveal}. However, this technique cannot distinguish two possible sources of preview effects: genuine pre-fixation integration of parafoveal information with the evolving sentence context, versus detection of a mismatch between the pre-saccadic preview and the post-saccadic foveal input \citep{schotter2015does, snell2017integration}. Furthermore, eye movement measures alone cannot reveal the neural dynamics of parafoveal processing or its timing within a fixation.

Here, we address these limitations using Rapid Invisible Frequency Tagging (RIFT) \citep{zhigalov2019rift, lozanosoldevilla2024rift}, a technique in which a visual stimulus is flickered at a frequency invisible to the observer (60~Hz) but detectable in the neural signal via MEG coherence analysis. By tagging the target word region with 60~Hz flicker and measuring the neural response during the pre-target fixation, we obtain a direct measure of the degree to which parafoveal attention is allocated to the upcoming word---before the reader fixates it.

\section{Results}

\subsection{Eye Movement Behaviour}

Standard eye-tracking measures confirmed the expected reading pattern (Table~\ref{tab:eyemove}).

\begin{table}[H]
\centering
\caption{Eye movement measures by condition (paired $t$-tests, $n = 34$).}
\label{tab:eyemove}
\begin{tabular}{@{}lcccc@{}}
\toprule
Measure & Congruent (ms) & Incongruent (ms) & $t(33)$ & $p$ \\
\midrule
First fixation on pre-target & $224.3$ & $228.1$ & $0.84$ & $0.407$ \\
First fixation on target & $241.7$ & $289.4$ & $5.99$ & $\mathbf{9.83 \times 10^{-7}}$ \\
Total gaze on target & $287.2$ & $371.8$ & $7.14$ & $\mathbf{2.91 \times 10^{-8}}$ \\
Refixation prob (pre-target) & $0.12$ & $0.23$ & $7.83$ & $\mathbf{5.04 \times 10^{-9}}$ \\
\bottomrule
\end{tabular}
\end{table}

No significant difference was found in first fixation duration on the pre-target word (Cohen's $d = 0.14$), but large effects emerged on the target word itself ($d = 1.03$) and in refixation probabilities ($d = 1.34$), confirming that incongruent words disrupted reading.

\subsection{RIFT Sensor Identification}

Sensors showing significant 60~Hz tagging responses were identified via cluster-based permutation comparing pre-target (flicker) and baseline (no-flicker) intervals ($p < 0.01$). Twenty-nine of 34 participants showed significant RIFT responses (Table~\ref{tab:rift_sensors}).

\begin{table}[H]
\centering
\caption{RIFT response characteristics.}
\label{tab:rift_sensors}
\begin{tabular}{@{}lc@{}}
\toprule
Parameter & Value \\
\midrule
Participants with RIFT response & 29 / 34 \\
Sensors per participant (mean $\pm$ SD) & $7.9 \pm 4.5$ \\
Sensor distribution & Posterior/occipital \\
Source location & Left BA18 (MNI: $[-9, -97, 3]$) \\
Source method & DICS beamforming \\
\bottomrule
\end{tabular}
\end{table}

\subsection{Pre-Target RIFT Congruency Effect}

During the pre-target fixation, 60~Hz coherence was significantly lower for incongruent than congruent upcoming target words (Table~\ref{tab:congruency}).

\begin{table}[H]
\centering
\caption{RIFT coherence during pre-target fixation by condition ($n = 29$).}
\label{tab:congruency}
\begin{tabular}{@{}lccc@{}}
\toprule
Condition & Mean Coherence ($\times 10^{-3}$) & SE & Statistical Test \\
\midrule
Congruent target & $4.72$ & $0.38$ & \multirow{2}{*}{$p = 0.011$ (cluster)} \\
Incongruent target & $3.91$ & $0.34$ & \\
\midrule
Onset of effect & \multicolumn{3}{c}{Within 100~ms of pre-target fixation onset} \\
\bottomrule
\end{tabular}
\end{table}

The congruency effect emerged within 100~ms of fixating the pre-target word, indicating very early semantic integration of parafoveal information with sentence context.

\subsection{Target Fixation RIFT Response}

No significant congruency difference was observed during direct fixation of the target word ($p > 0.15$), suggesting that the critical integration occurs during parafoveal preview rather than foveal processing.

\subsection{Position-Dependent Coherence Profile}

RIFT coherence at 60~Hz increased progressively as fixation position approached the target word (Table~\ref{tab:position}).

\begin{table}[H]
\centering
\caption{RIFT coherence by fixation position relative to the target word (N).}
\label{tab:position}
\begin{tabular}{@{}lcc@{}}
\toprule
Position & Coherence ($\times 10^{-3}$) & Significant Tagging? \\
\midrule
N$-$3 & $1.02$ & No \\
N$-$2 & $2.18$ & Marginal \\
N$-$1 (pre-target) & $4.31$ & \textbf{Yes} \\
N (target) & $3.87$ & Yes \\
N$+$1 & $2.43$ & Marginal \\
\bottomrule
\end{tabular}
\end{table}

\subsection{Individual Differences: RIFT and Reading Speed}

The magnitude of the RIFT congruency effect correlated with individual reading speed (Table~\ref{tab:individual}).

\begin{table}[H]
\centering
\caption{Correlation between RIFT congruency effect and reading speed ($n = 29$).}
\label{tab:individual}
\begin{tabular}{@{}lccc@{}}
\toprule
Measure & Spearman $\rho$ & $p$-value & Interpretation \\
\midrule
RIFT effect vs.\ reading speed & $-0.39$ & $\mathbf{0.037}$ & Larger effect $\rightarrow$ faster reading \\
\bottomrule
\end{tabular}
\end{table}

Participants with larger RIFT congruency effects (more negative coherence difference for incongruent) tended to read faster, linking parafoveal semantic integration to reading efficiency.

\subsection{Stimulus Properties and Counterbalancing}

Stimulus characteristics were carefully controlled across conditions (Table~\ref{tab:stimprops}).

\begin{table}[H]
\centering
\caption{Word characteristics for experimental stimuli (means across 80 target pairs).}
\label{tab:stimprops}
\begin{tabular}{@{}lcccc@{}}
\toprule
Property & Pre-target (Adj) & Target (Noun) & Cloze Prob & Freq (Zipf) \\
\midrule
Length (letters) & $5.8 \pm 1.2$ & $5.3 \pm 0.9$ & $0.04 \pm 0.03$ & $4.21 \pm 0.67$ \\
Congruent context & $5.7 \pm 1.1$ & $5.4 \pm 1.0$ & $0.05 \pm 0.04$ & $4.18 \pm 0.71$ \\
Incongruent context & $5.9 \pm 1.3$ & $5.2 \pm 0.8$ & $0.03 \pm 0.02$ & $4.24 \pm 0.63$ \\
$t$-test (Cong vs.\ Incong) & $t(79) = 0.84$ & $t(79) = 1.12$ & $t(79) = 0.67$ & $t(79) = 0.41$ \\
$p$-value & $0.40$ & $0.27$ & $0.51$ & $0.68$ \\
\bottomrule
\end{tabular}
\end{table}

No significant differences in word length, cloze probability, or frequency were found between congruent and incongruent conditions, confirming that the RIFT congruency effect reflects semantic integration rather than low-level stimulus confounds.

\section{Discussion}

This study provides the first direct neural evidence that semantic information from parafoveal words is integrated with sentence context before fixation during natural reading. Three key findings advance our understanding of parafoveal processing.

First, the reduced RIFT response for incongruent parafoveal words during the pre-target fixation demonstrates that semantic processing extends beyond the currently fixated word. Two mechanisms could explain this pattern: (1)~attentional redistribution, where semantically incongruent parafoveal content causes attention to shift back to the fixated word (covert regression), or (2)~an internal/external attention shift, where incongruent content demands more internal processing resources, reducing external visual responses \citep{schotter2012parafoveal}. Both accounts are consistent with pre-fixation semantic integration.

Second, the remarkably early onset of the congruency effect (within 100~ms) constrains the temporal dynamics of parafoveal semantic processing. This speed suggests that semantic integration with sentence context operates in parallel with ongoing foveal processing of the current word, consistent with models positing distributed attention across the perceptual span \citep{engbert2005swift, kliegl2006tracking}.

Third, the correlation between RIFT congruency effect magnitude and reading speed provides a mechanistic link between parafoveal processing and reading efficiency. Faster readers may extract and integrate more semantic information from the parafovea, allowing more efficient saccade planning and reduced fixation durations on upcoming words \citep{veldre2017parafoveal, reichle2009using}.

The absence of a congruency effect during direct fixation of the target word was unexpected but interpretable: once the reader fixates the target, the visual processing is dominated by foveal input and the RIFT tagging may no longer distinguish congruent from incongruent words because both receive full attentional resources.

\section{Methods}

\subsection{Participants}
Thirty-six native English speakers were recruited (34 analysed after exclusion; 23 female; mean age $22.5 \pm 2.8$ years; all right-handed; compensated at 15~GBP/hour).

\subsection{Stimuli}
One hundred sixty experimental sentences and 117 fillers were presented. Each sentence contained an unpredictable target noun (4--7 letters) preceded by an adjective (4--8 letters). Eighty target word pairs were counterbalanced across two sentence frames, producing congruent and incongruent conditions. Cloze probability was verified to be low for all targets.

\subsection{RIFT Paradigm}
Target word regions were flickered at 60~Hz using a 1440~Hz ProPixx projector. MEG was recorded with a 306-channel Elekta/MEGIN system. Eye movements were co-registered with MEG. Fixation cross duration was 1.2--1.6~s, followed by gaze-contingent box fixation (0.2~s) and sentence presentation.

\subsection{Analysis}
MEG coherence at 60~Hz between sensors and photodiode was computed using FieldTrip \citep{oostenveld2011fieldtrip}. Sensors with significant tagging were identified via cluster-based permutation \citep{maris2007nonparametric}. Source localisation used DICS beamforming \citep{gross2001dynamic}. Time-resolved coherence was computed over the pre-target fixation. Eye movement analyses used paired $t$-tests with Cohen's $d$ effect sizes.

\section{Conclusion}

Using RIFT combined with MEG and eye tracking during natural reading, we demonstrate that readers integrate semantic information from parafoveal words with sentence context before fixating those words, with effects emerging within 100~ms. The degree of parafoveal semantic integration predicts individual reading speed, establishing a functional link between this covert processing and reading efficiency.

\bibliographystyle{plainnat}
\bibliography{references}

\end{document}
